\clearpage
\thispagestyle{plain}
\phantomsection
\addcontentsline{toc}{chapter}{ABSTRACT}

{\centering\fontsize{12}{14}\selectfont\bfseries ABSTRACT\par}

\vspace{0.2cm}

% KONFIRMASI PENULIS: judul berbahasa Inggris adalah terjemahan \JudulPA, sesuaikan dengan judul resmi bila ada
{\noindent\raggedright
\textbf{IMPLEMENTATION OF A CONTINUOUS HEART RATE MONITORING SYSTEM BASED ON IOT EDGE COMPUTING FOR EARLY ARRHYTHMIA DETECTION}\par}

\vspace{0.2cm}

\noindent
\textbf{Name\hspace{1.55cm}: \NamaMahasiswa}\\
\textbf{NRP\hspace{1.75cm}: \NRP}\\
\textbf{Advisor\hspace{1.2cm}: \PembimbingSatu}\\
\textbf{\phantom{Advisor\hspace{1.2cm}: }\PembimbingDua}

\vspace{0.2cm}

Cardiovascular disease is the leading cause of death worldwide, and arrhythmia often occurs without symptoms, so continuous electrocardiogram (ECG) monitoring is needed for early detection. Existing monitoring systems are generally cloud-centric, with high latency and a dependence on internet connectivity. Arrhythmia detection models are also rarely tested under microcontroller constraints, and edge-to-cloud data pipelines are seldom tolerant of network disruption. This study builds a single-lead heart monitoring system based on edge computing using a TinyML approach. ECG preprocessing and arrhythmia classification run on-device on an ESP32-S3, and the results are sent through an MQTT pipeline designed to tolerate connection loss. A Butterworth filter and Pan-Tompkins R-peak detection run on the microcontroller, and a hybrid CNN with a morphology branch and a rhythm branch is quantized to INT8. Beats that have not been delivered are held in a local buffer in PSRAM. The float32 model on the MIT-BIH DS2 set, which contains 22 patients unseen during training, achieves a recall of 0.6996, an F1 of 0.6594, and an AUC of 0.9334. INT8 quantization yields a 22.94~KB model with a recall difference of $+$0.0108, and on-device inference takes 26.6~ms per beat. Preprocessing runs on the ESP32-S3 at 3.0\% of the time budget with no lost samples. All four official body-signal recordings from one recording session pass the signal quality index in every 10-second window, and the electrode-contact calibration gate opens monitoring about 60 seconds after good contact. The preprocessing and model objectives are achieved, while the MQTT pipeline has been designed and its test results will be reported in the next stage.
% sumber: terjemahan bab/abstrak.tex, angka identik dengan abstrak berbahasa Indonesia dan Bab 4 (desimal memakai titik sesuai kaidah bahasa Inggris)

\vspace{0.25cm}
\noindent
\textbf{Keywords:} ECG, IoT, Edge Computing, TinyML, ESP32-S3, Arrhythmia, MQTT

\clearpage
